% =============================================================================
% Report.tex -- Madison 2026 Assessment Methodology Fact-Check
%
% Compile from the paper/ directory with:
%   pdflatex Report.tex
%   pdflatex Report.tex   (second pass resolves cross-references)
%
% Numbers are loaded from ../outputs/tables/headlines.tex (auto-generated by
% scripts/03_build_report_assets.py). Run all three numbered scripts first.
% =============================================================================
\documentclass[11pt, letterpaper]{article}

% --- Layout ------------------------------------------------------------------
\usepackage[top=1.2in, bottom=1.2in, left=1.2in, right=1.2in]{geometry}
\usepackage{setspace}
\setstretch{1.15}
\usepackage[activate=false]{microtype}
\usepackage{parskip}
\setlength{\parskip}{6pt}

% --- Overflow prevention -----------------------------------------------------
\usepackage{xurl}
\usepackage[htt]{hyphenat}
\sloppy
\emergencystretch=3em

% --- Tables ------------------------------------------------------------------
\usepackage{booktabs}
\usepackage{tabularx}
\usepackage{array}

% --- Typography / colour -----------------------------------------------------
\usepackage{xcolor}
\definecolor{madisonblue}{HTML}{003A70}
\definecolor{madisonred}{HTML}{B71C1C}
\definecolor{confirmed}{HTML}{1B5E20}

% --- Hyperlinks --------------------------------------------------------------
\usepackage[
  colorlinks=true,
  linkcolor=madisonblue,
  citecolor=madisonblue,
  urlcolor=madisonblue,
  pdftitle={Madison 2026 Assessment Methodology Change: A Data Fact-Check},
  pdfauthor={Russell Richie}
]{hyperref}

% --- Section formatting ------------------------------------------------------
\usepackage{titlesec}
\titleformat{\section}{\large\bfseries\color{madisonblue}}{\thesection.}{0.5em}{}
\titleformat{\subsection}{\normalsize\bfseries}{\thesubsection.}{0.5em}{}

% --- Verdict box command -----------------------------------------------------
\newcommand{\verdict}[2]{%
  \noindent\textbf{Claim:} \textit{#1}\\[2pt]
  \noindent\textbf{Finding:} #2\par}

% --- Load generated number macros --------------------------------------------
% Macro naming convention: all letters only (LaTeX forbids digits in \commands).
%   Years: TwentyFour / TwentyFive / TwentySix
%   YoY: Prev (2024->2025) / Curr (2025->2026)
%   Examples: ordinals One/Two/Three, years Prev/Curr
\InputIfFileExists{../outputs/tables/headlines.tex}{}{%
  \typeout{WARNING: headlines.tex not found. Run scripts 01-03 first.}}

% =============================================================================
\begin{document}

% --- Title -------------------------------------------------------------------
\begin{center}
  {\LARGE\bfseries Madison 2026 Assessment Methodology Change}\\[0.4em]
  {\large\bfseries A Data Fact-Check}\\[1em]
  {\normalsize Russell Richie \quad\textbullet\quad May 2026}\\[0.5em]
  {\small\itshape
    Data: City of Madison Office of the City Assessor, Property Tax Base
    Reports 2024--2025; City of Madison Open Data Portal, Tax Parcels
    (Assessor Property Information), ArcGIS Feature Service
    \url{https://maps.cityofmadison.com/arcgis/rest/services/Public/OPEN_DATA2/FeatureServer/0}}
\end{center}

\bigskip
\noindent\rule{\textwidth}{0.8pt}
\bigskip

% =============================================================================
\section{The Claim}
% =============================================================================

A May 2026 letter to the editor in the \textit{Capital Times} (Ginny White,
``Assessments Designed to Build at Any Cost'') alleged that Madison's City
Assessor had quietly changed assessment methodology:

\begin{quote}
``\ldots the city quietly changed how it assesses property taxes. It kept
the total assessment the same as last year but raised the land assessment
and lowered the improvement assessment by equal amounts. This `land value
tax' approach is designed to promote development, penalize land speculation
and capture rising land value.''
\end{quote}

We evaluate each specific claim against publicly available assessment data.

% =============================================================================
\section{Data Sources}
% =============================================================================

Two complementary sources were used:

\textbf{Annual Property Tax Base Reports.}
The City Assessor publishes a summary report each May.
The 2025 report (dated April 19, 2024) and 2026 report (dated May 19, 2025)
provide city-wide totals by property class but do not disaggregate assessments
into land versus improvement components.

\textbf{Open Data Portal --- Tax Parcel Feature Service.}
The City of Madison publishes a live ArcGIS feature layer containing all
82,249 taxable parcels. Each parcel record includes three years of land
and improvement values in separate fields:
\texttt{CurrentLand}/\texttt{CurrentImpr}/\texttt{CurrentTotal} (2026 assessment),
\texttt{PreviousLand}/\texttt{PreviousImpr}/\texttt{PreviousTotal} (2025),
and \texttt{PreviousLand2}/\texttt{PreviousImpr2}/\texttt{PreviousTotal2} (2024).
Aggregate statistics and counts were queried directly via the
ArcGIS REST API \texttt{/query} endpoint.

% =============================================================================
\section{Findings}
% =============================================================================

% -----------------------------------------------------------------------------
\subsection{City-Wide Assessment Trends}
% -----------------------------------------------------------------------------

Table~\ref{tab:city-wide} reports city-wide sums of land and improvement
values across three assessment years derived from the open data portal.
The totals are consistent with the official Property Tax Base Reports
(the 2025 parcel sum of \TotalTwentyFive{} matches the published \$49.0B figure).

\begin{table}[htbp]
  \centering
  \caption{City-wide assessed value, Madison, WI --- land vs.\ improvements
           (all 82,249 taxable parcels). Source: City of Madison Open Data
           Portal, Tax Parcels Feature Service, queried May 2026.}
  \label{tab:city-wide}
  \begin{tabularx}{\textwidth}{@{} l X X X @{}}
    \toprule
    & \textbf{2024} & \textbf{2025} & \textbf{2026} \\
    \midrule
    Land                & \LandTwentyFour{} & \LandTwentyFive{} & \LandTwentySix{} \\
    \quad Year-over-year change & --- & \LandYoYPrev{} & $\mathbf{\LandYoYCurr{}}$ \\[4pt]
    Improvements        & \ImprTwentyFour{} & \ImprTwentyFive{} & \ImprTwentySix{} \\
    \quad Year-over-year change & --- & \ImprYoYPrev{} & $\mathbf{\ImprYoYCurr{}}$ \\[4pt]
    \textbf{Total}      & \TotalTwentyFour{} & \TotalTwentyFive{} & \TotalTwentySix{} \\
    \quad Year-over-year change & --- & \TotalYoYPrev{} & \TotalYoYCurr{} \\[4pt]
    Land as \% of total & \LandPctTwentyFour{} & \LandPctTwentyFive{} & \textbf{\LandPctTwentySix{}} \\
    \bottomrule
  \end{tabularx}
  \smallskip\noindent\parbox{\linewidth}{\small
  The pattern reverses sharply between 2025 and 2026: improvements drove
  value growth in 2025 (\ImprYoYPrev{} vs.\ \LandYoYPrev{} for land); in 2026,
  land grew far faster than improvements (\LandYoYCurr{} vs.\ \ImprYoYCurr{}).}
\end{table}

The most recent assessment cycle (2026) exhibits a dramatic reversal:
land values jumped \LandYoYCurr{} city-wide while improvement values barely moved
(\ImprYoYCurr{}), compared to the prior year when the ratio was essentially
the opposite. The land share of total assessed value rose from \LandPctTwentyFive{}
to \LandPctTwentySix{}.

% -----------------------------------------------------------------------------
\subsection{Parcel-Level Pattern}
% -----------------------------------------------------------------------------

The open data portal enables direct counting of parcels matching
the specific pattern the letter describes: improvements decrease,
land increases, and total stays approximately or exactly flat.
Table~\ref{tab:parcel-counts} reports these counts for the two most
recent year-over-year comparisons.

\begin{table}[htbp]
  \centering
  \caption{Parcels matching each component of the letter's claim,
           2024$\to$2025 vs.\ 2025$\to$2026.
           Source: City of Madison Tax Parcels Feature Service, REST API
           aggregate queries.}
  \label{tab:parcel-counts}
  \begin{tabularx}{\textwidth}{@{} X r r @{}}
    \toprule
    Criterion & 2024$\to$2025 & 2025$\to$2026 \\
    \midrule
    Improvements $\downarrow$ \textbf{and} land $\uparrow$ & \PrevBoth{} & \textbf{\CurrBoth{}} \\
    Impr.\ $\downarrow$, land $\uparrow$, total within $\pm3\%$ & --- & \CurrNearFlat{} \\
    Impr.\ $\downarrow$, land $\uparrow$, total \textbf{exactly} unchanged & \textbf{\PrevExact{}} & \textbf{\CurrExact{}} \\
    \midrule
    Total parcels & 82,249 & 82,249 \\
    \bottomrule
  \end{tabularx}
  \smallskip\noindent\parbox{\linewidth}{\small
  The salient comparison is the jump from \PrevBoth{} to \CurrBoth{} and
  from \PrevExact{} to \CurrExact{} in the exactly-flat category.}
\end{table}

The scale of the change is striking. In the 2024$\to$2025 cycle, just
\textbf{\PrevBoth{} parcels} had improvements fall while land rose; in
2025$\to$2026, \textbf{\CurrBoth{} parcels} --- \CurrBothPct{} of all taxable
parcels --- show this pattern, a $\BothMultiplier{}$ increase. Of those,
\textbf{\CurrExact{} parcels} had their improvement value lowered and land
value raised by \emph{exactly} equal and opposite amounts, leaving the total
assessment unchanged to the dollar.

% -----------------------------------------------------------------------------
\subsection{Example Parcels}
% -----------------------------------------------------------------------------

Table~\ref{tab:examples} shows three representative parcels drawn from the
live feature service. In each case the shift from improvements to land is
exact: the dollar increase in land equals the dollar decrease in improvements,
and the total assessment is unchanged.

\begin{table}[htbp]
  \centering
  \caption{Three representative parcels where the 2026 assessment held the
           total exactly flat while shifting value from improvements to land.
           Source: City of Madison Tax Parcels Feature Service.}
  \label{tab:examples}
  \begin{tabularx}{\textwidth}{@{} l X X X X @{}}
    \toprule
    Address & Land & Impr. & Total & Shift \\
    \midrule
    \multicolumn{5}{@{}l}{\textit{2025 (Previous) assessment}} \\[2pt]
    \ExAddrOne{}   & \ExLandPrevOne{}   & \ExImprPrevOne{}   & \ExTotalPrevOne{}   & --- \\
    \ExAddrTwo{}   & \ExLandPrevTwo{}   & \ExImprPrevTwo{}   & \ExTotalPrevTwo{}   & --- \\
    \ExAddrThree{} & \ExLandPrevThree{} & \ExImprPrevThree{} & \ExTotalPrevThree{} & --- \\[4pt]
    \multicolumn{5}{@{}l}{\textit{2026 (Current) assessment}} \\[2pt]
    \ExAddrOne{}   & \ExLandCurrOne{}   & \ExImprCurrOne{}   & \ExTotalCurrOne{}   & $+$\ExShiftOne{} land; $-$\ExShiftOne{} impr. \\
    \ExAddrTwo{}   & \ExLandCurrTwo{}   & \ExImprCurrTwo{}   & \ExTotalCurrTwo{}   & $+$\ExShiftTwo{} land; $-$\ExShiftTwo{} impr. \\
    \ExAddrThree{} & \ExLandCurrThree{} & \ExImprCurrThree{} & \ExTotalCurrThree{} & $+$\ExShiftThree{} land; $-$\ExShiftThree{} impr. \\
    \bottomrule
  \end{tabularx}
  \smallskip\noindent\parbox{\linewidth}{\small
  Each parcel shows an exact dollar-for-dollar shift from improvements
  to land, leaving the total assessment unchanged.}
\end{table}

% -----------------------------------------------------------------------------
\subsection{Breakdown by Property Class and Use}
% -----------------------------------------------------------------------------

The methodology shift is concentrated almost entirely in the
\textbf{Residential} class. Table~\ref{tab:by-class} shows land and
improvement assessment changes by the three taxable property classes with
meaningful 2026 values (Exempt and Manufacturing parcels carry zeroed-out
current values and are excluded).

\begin{table}[htbp]
  \centering
  \caption{Assessment changes by property class, 2025$\to$2026.
           Exempt and Manufacturing parcels excluded (2026 values are zero
           for fully tax-exempt properties).
           Source: City of Madison Tax Parcels Feature Service.}
  \label{tab:by-class}
  \begin{tabularx}{\textwidth}{@{} X r r r r r @{}}
    \toprule
    Class (parcels) & Land YoY & Impr.\ YoY & Total YoY & Impr.$\downarrow$+Land$\uparrow$ & Exact flat \\
    \midrule
    Residential (\ParcelsRes{})  & \LandYoYRes{}  & \ImprYoYRes{}  & \TotalYoYRes{}  & \BothRes{} (\BothPctRes{})  & \ExactRes{} \\
    Commercial  (\ParcelsComm{}) & \LandYoYComm{} & \ImprYoYComm{} & \TotalYoYComm{} & \BothComm{} (\BothPctComm{}) & \ExactComm{} \\
    Agriculture (\ParcelsAg{})   & \LandYoYAg{}   & \ImprYoYAg{}   & \TotalYoYAg{}   & \BothAg{} (\BothPctAg{})    & \ExactAg{} \\
    \bottomrule
  \end{tabularx}
  \smallskip\noindent\parbox{\linewidth}{\small
  Exempt and Manufacturing parcels carry zeroed-out 2026 values (fully tax-exempt)
  and are excluded. The land-up/improvement-down pattern is overwhelmingly residential:
  only \BothComm{} of \ParcelsComm{} commercial parcels (\BothPctComm{}) show the
  pattern, versus \BothRes{} of \ParcelsRes{} residential (\BothPctRes{}).}
\end{table}

Table~\ref{tab:by-use} further disaggregates residential parcels by
property use. The effect is concentrated in \textbf{single-family homes}
and \textbf{small multi-family} (2- and 3-unit) properties.
Condominiums and vacant lots are largely unaffected.

\begin{table}[htbp]
  \centering
  \caption{Assessment changes within the Residential class by property use,
           2025$\to$2026. Source: City of Madison Tax Parcels Feature Service.}
  \label{tab:by-use}
  \begin{tabularx}{\textwidth}{@{} X r r r r r @{}}
    \toprule
    Property use (parcels) & Land YoY & Impr.\ YoY & Total YoY & Impr.$\downarrow$+Land$\uparrow$ & Exact flat \\
    \midrule
    Single family (\ParcelsSF{})    & \LandYoYSF{}        & \ImprYoYSF{}        & \TotalYoYSF{}        & \BothSF{} (\BothPctSF{})               & \ExactSF{} \\
    Condominium   (\ParcelsCondo{}) & \LandYoYCondo{}     & \ImprYoYCondo{}     & \TotalYoYCondo{}     & \BothCondo{} (\BothPctCondo{})         & \ExactCondo{} \\
    2-Unit  (\ParcelsTwoUnit{})     & \LandYoYTwoUnit{}   & \ImprYoYTwoUnit{}   & \TotalYoYTwoUnit{}   & \BothTwoUnit{} (\BothPctTwoUnit{})     & \ExactTwoUnit{} \\
    3-Unit  (\ParcelsThreeUnit{})   & \LandYoYThreeUnit{} & \ImprYoYThreeUnit{} & \TotalYoYThreeUnit{} & \BothThreeUnit{} (\BothPctThreeUnit{}) & \ExactThreeUnit{} \\
    Vacant  (\ParcelsVac{})         & \LandYoYVac{}       & \ImprYoYVac{}       & \TotalYoYVac{}       & \BothVac{} (\BothPctVac{})             & \ExactVac{} \\
    \bottomrule
  \end{tabularx}
  \smallskip\noindent\parbox{\linewidth}{\small
  Single-family (\BothPctSF{}) and small multi-family (\BothPctTwoUnit{}--\BothPctThreeUnit{})
  properties show the pattern at high rates. Condominiums (\BothPctCondo{}) and vacant
  lots (\BothPctVac{}) are effectively excluded from the shift. Among condominiums
  that do show the pattern (\BothCondo{} parcels), nearly all (\ExactCondo{}) are
  exact-flat shifts, suggesting a distinct assessment mechanism for that property type.}
\end{table}

% =============================================================================
\section{Verdict}
% =============================================================================

\bigskip
\verdict{%
  ``It kept the total assessment the same as last year.''}{%
  \textbf{\color{madisonred}Partially true, city-wide false.}
  City-wide assessed value rose \TotalYoYCurr{} (\TotalTwentyFive{} $\to$ \TotalTwentySix{}). However,
  for \textbf{\CurrExact{} individual parcels} the total assessment was held
  exactly flat while land and improvement values were shifted by equal
  and opposite amounts. The author was accurately describing their
  own parcel's experience; the characterization does not extend to
  the city as a whole.}

\bigskip
\verdict{%
  ``Raised the land assessment and lowered the improvement assessment
  by equal amounts.''}{%
  \textbf{\color{confirmed}Confirmed as a real methodology change.}
  \CurrBoth{} parcels (\CurrBothPct{} of all taxable parcels) had improvements
  decrease and land increase in the 2026 assessment cycle, versus just
  \PrevBoth{} parcels in 2025 --- a $\BothMultiplier{}$ increase. For \CurrExact{}
  of those parcels the shift was exact to the dollar. This is clearly a
  deliberate and systematic reassessment, not a statistical artifact.}

\bigskip
\verdict{%
  ``This `land value tax' approach is designed to promote development
  and penalize land speculation.''}{%
  \textbf{\color{madisonblue}Characterization is reasonable.}
  The observed pattern --- land taxed more heavily, improvements
  less so --- is structurally equivalent to a partial land value tax
  shift applied at the sub-parcel level. The most plausible driver is
  the 2023 Transit-Oriented Development (TOD) overlay zoning along
  BRT corridors, which increased permitted density and hence land
  market value in affected neighborhoods. Wisconsin municipalities
  cannot adopt a true LVT independently, but a reassessment methodology
  that captures rising land values while holding improvement assessments
  flat achieves similar incentive effects.}

\bigskip\noindent\rule{\textwidth}{0.4pt}
\smallskip
\noindent{\small\itshape
Data accessed May 2026. The open data feature service reflects the 2026
assessment cycle (assessed as of January 1, 2026; notices mailed May 2026).
All API queries are reproducible against the public endpoint at
\url{https://maps.cityofmadison.com/arcgis/rest/services/Public/OPEN_DATA2/FeatureServer/0}.}

\end{document}
